% Options for packages loaded elsewhere
\PassOptionsToPackage{unicode}{hyperref}
\PassOptionsToPackage{hyphens}{url}
\documentclass[
]{article}
\usepackage{xcolor}
\usepackage[margin=25mm]{geometry}
\usepackage{amsmath,amssymb}
\setcounter{secnumdepth}{-\maxdimen} % remove section numbering
\usepackage{iftex}
\ifPDFTeX
  \usepackage[T1]{fontenc}
  \usepackage[utf8]{inputenc}
  \usepackage{textcomp} % provide euro and other symbols
\else % if luatex or xetex
  \usepackage{unicode-math} % this also loads fontspec
  \defaultfontfeatures{Scale=MatchLowercase}
  \defaultfontfeatures[\rmfamily]{Ligatures=TeX,Scale=1}
\fi
\usepackage{lmodern}
\ifPDFTeX\else
  % xetex/luatex font selection
\fi
% Use upquote if available, for straight quotes in verbatim environments
\IfFileExists{upquote.sty}{\usepackage{upquote}}{}
\IfFileExists{microtype.sty}{% use microtype if available
  \usepackage[]{microtype}
  \UseMicrotypeSet[protrusion]{basicmath} % disable protrusion for tt fonts
}{}
\makeatletter
\@ifundefined{KOMAClassName}{% if non-KOMA class
  \IfFileExists{parskip.sty}{%
    \usepackage{parskip}
  }{% else
    \setlength{\parindent}{0pt}
    \setlength{\parskip}{6pt plus 2pt minus 1pt}}
}{% if KOMA class
  \KOMAoptions{parskip=half}}
\makeatother
\usepackage{longtable,booktabs,array}
\newcounter{none} % for unnumbered tables
\usepackage{calc} % for calculating minipage widths
% Correct order of tables after \paragraph or \subparagraph
\usepackage{etoolbox}
\makeatletter
\patchcmd\longtable{\par}{\if@noskipsec\mbox{}\fi\par}{}{}
\makeatother
% Allow footnotes in longtable head/foot
\IfFileExists{footnotehyper.sty}{\usepackage{footnotehyper}}{\usepackage{footnote}}
\makesavenoteenv{longtable}
\usepackage{graphicx}
\makeatletter
\newsavebox\pandoc@box
\newcommand*\pandocbounded[1]{% scales image to fit in text height/width
  \sbox\pandoc@box{#1}%
  \Gscale@div\@tempa{\textheight}{\dimexpr\ht\pandoc@box+\dp\pandoc@box\relax}%
  \Gscale@div\@tempb{\linewidth}{\wd\pandoc@box}%
  \ifdim\@tempb\p@<\@tempa\p@\let\@tempa\@tempb\fi% select the smaller of both
  \ifdim\@tempa\p@<\p@\scalebox{\@tempa}{\usebox\pandoc@box}%
  \else\usebox{\pandoc@box}%
  \fi%
}
% Set default figure placement to htbp
\def\fps@figure{htbp}
\makeatother
\setlength{\emergencystretch}{3em} % prevent overfull lines
\providecommand{\tightlist}{%
  \setlength{\itemsep}{0pt}\setlength{\parskip}{0pt}}
\usepackage{bookmark}
\IfFileExists{xurl.sty}{\usepackage{xurl}}{} % add URL line breaks if available
\urlstyle{same}
\hypersetup{
  hidelinks,
  pdfcreator={LaTeX via pandoc}}

\author{}
\date{}

\begin{document}

\section{The Periodic Table of Waves: A Hypothesis for Particle
Classification by Winding Addresses and Observation
Clocks}\label{the-periodic-table-of-waves-a-hypothesis-for-particle-classification-by-winding-addresses-and-observation-clocks}

\textbf{Author}: Noriaki Kihara \textbf{Date}: August 6, 2026
\textbf{Type}: Hypothesis-proposal paper (supported by numerical
experiments; not a proof paper) \textbf{Version DOI}:
10.5281/zenodo.21822359 \textbf{Concept DOI}: 10.5281/zenodo.21822358

\begin{center}\rule{0.5\linewidth}{0.5pt}\end{center}

\subsection{Abstract}\label{abstract}

By what are particles classified? From numerical experiments on wave
dynamics that assume no particle concept whatsoever (two published
systems, used as-is), this paper proposes the \textbf{``periodic table
of waves''} --- a hypothesis for particle classification by winding
addresses and observation clocks. The starting point was a sign-vector
thought experiment {[}1{]} classifying particles by which of the six
axes (x, y, z, t, R, Q) carries their winding. This paper rebuilds that
classification from wave forms alone, by measurement, arriving at five
pillars. \textbf{Pillar 1 (periodic law)}: the foundation of
classification is a measured universal clock --- the collective clock
\(\omega\) = \(\pi\)/72 per step holds across the entire range N = 4 to
144 (long-window \(\pm\)\textasciitilde0.1\%) --- and species are
classified by rational addresses on a single U\^{}144 clock lattice and
their relation to the clock. \textbf{Pillar 2 (charge rule)}: charge is
Q\_em = m/3 (m the raw winding number), read by the dominant observation
clock (denominator 3 --- based on the previously measured tongue-width
ratio of 461) as ``3 windings = 1 elementary charge.'' It follows that
the u-type m = +2 reads +2/3, the d-type m = -1 reads -1/3, the electron
m = -3 reads -1, and \textbf{confinement and fractional charge become
two faces of the single fact of ``mod-3 readability''} (the bookkeeping
checks for the proton uud = +3 \(\to\) +1 and the neutron udd = 0 pass
automatically). In the numerical experiments, quark-type addresses
remain exactly unreadable in isolation (readable fraction 0.0000
persists), while in the sea they become readable only in integer charge
units, to the extent that hadronization (formation of m \(\equiv\) 0
composites) proceeds. \textbf{Pillar 3 (statistics rule)}: the
Fermi/Bose distinction lives not in spatial rotation but in the
\textbf{double cover of the clock}. The state phase is quantized to \{0,
\(\pi\)\} at every observational recurrence (measured to 3 digits);
fermionic species (odd harmonics) visit both sheets (covering degree 2),
bosonic species do not (degree 1). A pure m = 0 fermionic-band species
also shows covering degree 2, establishing the row of the neutral
fermion (neutrino). \textbf{Pillar 4 (mass rule)}: mass is not an
attribute of the address but a relational quantity with the sea (vacuum)
--- isolated addresses are all unitarily equivalent under winding shifts
(measured), and in-sea properties depend only on the divisor class
gcd(m, 16) (odd addresses agree to 5 digits). \textbf{Pillar 5 (the
table)}: the table is a double signboard --- an isolated stable-element
table and an in-sea lifetime table --- on which a 62-species
Standard-Model assignment is proposed. Under this hypothesis, the
problem left open by the previous paper {[}5{]} --- ``charge does not
stabilize at \(\pm\)\textasciitilde1'' --- turns from a contradiction
into a prediction, and the 1/3-charge problem of quarks connects to
confinement through a single mechanism. This paper claims no proof. What
it claims is that seven independently measured structures connect, under
a single classification hypothesis, simultaneously to the particle
structure predicted by the Standard Model, to the open problem of the
previous paper, and to the sign-vector thought experiment. All
experiments are published as deterministic scripts, and the entire
process --- including 13 refutations --- is reproducible. The closing
section fixes the requirements that the universal interaction function
of the next paper must satisfy in order to verify this classification
hypothesis dynamically.

\begin{center}\rule{0.5\linewidth}{0.5pt}\end{center}

\subsection{1. Introduction --- why a ``periodic table,'' and why a
hypothesis
proposal}\label{introduction-why-a-periodic-table-and-why-a-hypothesis-proposal}

\subsubsection{1.1 Starting point: the sign-vector thought
experiment}\label{starting-point-the-sign-vector-thought-experiment}

A preceding thought experiment {[}1{]} proposed classifying particles by
which of the six axes (x, y, z, t, R, Q) carries their winding ---
photon \(\gamma\) = (0,0,0,0,R,0), W\(\pm\) = (0,0,0,\(\pm\)t,0,0),
gluon g = (0,0,0,0,0,Q), graviton G = (0,0,0,\(\pm\)t,\(\pm\)R,0). The
classification is attractive but remained at the level of a thought
experiment. Can particle species be classified as \textbf{measurements}
of wave dynamics? That is the question of this paper.

\subsubsection{1.2 Central proposition}\label{central-proposition}

All pillars of this paper are bound by one proposition:

\begin{quote}
\textbf{Particle species are classified not as fixed attributes of the
waves themselves, but by how an observation clock reads winding states.}
Charge, confinement, statistics, mass, and lifetime are not independent
labels; they decompose into \textbf{state-side address \(\times\)
observation clock \(\times\) relation to the sea}.
\end{quote}

Charge quantization exists neither in the state space alone nor in the
measuring apparatus alone --- it exists in \textbf{the relation between
the cyclic structure of the state and the divisor structure of the
clock} (§4). Statistics lives in the double cover of the clock, not in
spatial rotation (§5), and mass is differentiated by the sea (§6). This
is the most important theoretical proposal of this paper.

\subsubsection{1.3 Methodological declaration: no particles
assumed}\label{methodological-declaration-no-particles-assumed}

This paper introduces no new dynamics. It uses two published systems ---
the N-body relational-wave dynamics {[}2{]} (used for measurements on
the three-directional condensate) and the closed-form exact solution of
the two-body universal inelastic map {[}3{]} (used for all charge,
statistics, and confinement experiments) --- both built on the anonymity
principle: ``wave forms are the only input, and there is no
special-casing inside.'' The concepts of particle, charge, and spin do
not enter the dynamics. We measure how they arise \textbf{as readouts}.

\subsubsection{1.4 Declaration of character, and historical
lessons}\label{declaration-of-character-and-historical-lessons}

This is a \textbf{hypothesis-proposal paper}. Attempts to classify
particles have a history of both success and failure. Gell-Mann's
Eightfold Way {[}9{]} succeeded, predicting the \(\Omega^-\) from SU(3)
classification. Kelvin's vortex atoms {[}8{]} --- atoms as knots of
vortices, a topological classification --- were beautiful, but the
dynamics was wrong. Preon models {[}14{]} attempted 62-species-scale
full assignments and faded. \textbf{A beautiful classification does not
guarantee correct dynamics.} That this paper remains a hypothesis
proposal rather than a proof, strictly separates readout-level from
dynamics-level claims (§2.3), and states explicit falsification
conditions (§11), is due to this lesson. The value of the hypothesis
should be measured by the breadth of its connections and the clarity of
its falsification conditions.

\subsection{2. Methods and conventions}\label{methods-and-conventions}

\subsubsection{2.1 Dynamics (read-only)}\label{dynamics-read-only}

\begin{itemize}
\tightlist
\item
  \textbf{N-body system} {[}2{]}: closure dynamics of N relational waves
  (edges of the complete graph), forming a three-directional metastable
  condensate. Used for the clock-universality scan (§3).
\item
  \textbf{Two-body system} {[}3{]}: the closed-form exact solution of
  the pointwise interaction \(\delta a\) = -2R·Im(\(\bar b a\))·b of a
  two-channel wave (a, b). States live on a (\(\chi\), \(\eta\)) grid,
  and the \(\eta\) winding number m is the object of charge readout.
  Used for all experiments of §4--6.
\end{itemize}

Both are deterministic; every script and result can be re-executed
standalone according to the attached inventory (the experiment inventory
file (jikken\_ichiran\_v1.md)).

\subsubsection{2.2 Machine-verified
conventions}\label{machine-verified-conventions}

To avoid convention collisions, two points were machine-checked
(attachment \#25): (i) bundle-construction bin numbers and \(\chi\)
frequencies correspond with a \(\pm\)\textasciitilde1 shift and inverted
parity (even-bin bundles \(\to\) odd frequencies and vice versa, power
fraction 1.0000). The text standardizes on the
\textbf{\(\chi\)-frequency representation}. Fermionic = even \(\chi\)
frequency (\(\ge\)\textasciitilde4) = odd harmonics {[}7{]}. (ii)
Bundles carry exactly one unit of \(\eta\) winding m = +1 by
construction (\(\eta\) occupancy (1, 1.000)). A species of arbitrary
winding m is constructed by the winding shift e\^{}\{i(m-1)\(\eta\)\},
not by projection.

\subsubsection{2.3 The most important convention: readout level
vs.~dynamics
level}\label{the-most-important-convention-readout-level-vs.-dynamics-level}

The claims of this paper are restricted to the \textbf{readout level}
--- ``particle classification arises as a relation between observation
clock and address.'' The \textbf{dynamics level} --- ``62 species are
generated, decay, and interact as the spectrum of a single dynamics''
--- is outside the scope of this paper and is delegated to the next
paper (the universal interaction function). The closing section (§13)
fixes its requirement table on the basis of measurements.

\subsection{3. Pillar 1: the foundation of the periodic law --- the
universal clock
(measured)}\label{pillar-1-the-foundation-of-the-periodic-law-the-universal-clock-measured}

The collective clock of the N-body condensate was scanned from N = 4 to
144 (28 points).

\textbf{Measurement}: \(\omega_{\rm clock}\) = \(\pi\)/72 per step holds
across the entire range --- ratio 0.988 \(\pm\) 0.028 in short windows
(T = 4000), and 0.9992--1.0013 (\textbf{\(\pm\)\textasciitilde0.1\%}) in
long windows (T = 42000). One clock revolution = 144 steps, independent
of the effective energy (N).

\begin{figure}
\centering
\pandocbounded{\includegraphics[keepaspectratio,alt={Figure P1}]{figs/fig_p1_clock_universality_v1.png}}
\caption{Figure P1}
\end{figure}

\textbf{Figure P1}: Clock universality: \(\omega\)/(\(\pi\)/72) = 1 for
all N.

The long-window census further showed that in the low-energy sea (N
\(\le\) 16) exactly \textbf{one} species is long-time stationary: the
massless ground species locked to the clock, \(\rho\) = 1/1
(\textbar{}\(\rho\)-1\textbar{} \textless{} 9\(\times\)\(~10^{-4}\);
mass degree \(~10^{-11}\)). The massive sidebands seen in short windows
are all finite-lifetime transients, entrained into the clock (Figure
P2). The refutation process (short-window apparent addresses and mass
laws vanishing in long windows) is fully recorded in §12.

\begin{figure}
\centering
\pandocbounded{\includegraphics[keepaspectratio,alt={Figure P2}]{figs/fig_p2_stable_vs_resonance_v1.png}}
\caption{Figure P2}
\end{figure}

\textbf{Figure P2}: Short-window sidebands (resonances) collapse to 1/1
in long windows.

\textbf{Hypothesis (Pillar 1)}: the classifier of particle species is
the rational address of U\^{}n = I, and the foundation is the single
U\^{}144 clock lattice. Species are classified by ``address (winding)
\(\times\) relation to the clock.''

\subsection{4. Pillar 2: the charge rule --- Q\_em = m/3 and confinement
= mod-3
readability}\label{pillar-2-the-charge-rule-q_em-m3-and-confinement-mod-3-readability}

\subsubsection{4.1 Multivaluedness of charge (state-side
measurement)}\label{multivaluedness-of-charge-state-side-measurement}

The charged species (winding q = +1) is longer-lived than neutral
transients by orders of magnitude (\(\tau\) \(\approx\)
1.3\(\times\)\(~10^{4}\) collisions; 85\% retention), but not permanent.
Its decay is not disappearance but transport to the partner (+3) by the
walk of the sum rule m* = 2m\_B - m\_s {[}4{]} (Figure P3). Moreover,
isolated pure winding species \textbf{replicate themselves exactly for
any m} (retention 1.000; \(\tau\) numerically infinite), and even the
sea-free mixture \{+1, +3\} does not leak --- \textbf{the walk is driven
only by the m = 0 sea} (leak destinations \{0, -1,
+2\}/\{\(\pm\)\textasciitilde4\} agree completely with the sum-rule
products). The previous paper's {[}5{]} §9.5 finding --- ``charge does
not stabilize at \(\pm\)\textasciitilde1'' --- is a direct consequence
of this sea-driven walk.

\begin{figure}
\centering
\pandocbounded{\includegraphics[keepaspectratio,alt={Figure P3}]{figs/fig_p3_charged_lifetime_walk_v1.png}}
\caption{Figure P3}
\end{figure}

\textbf{Figure P3}: Quasi-stability of the charged species and the
sum-rule walk (+1 \(\to\) +3).

\subsubsection{4.2 The conservation law of winding charge is
cyclic}\label{the-conservation-law-of-winding-charge-is-cyclic}

The true conservation law of winding charge is not integer but
\textbf{cyclic, mod ne (the \(\eta\) register order, 16 in this
system)}. The pointwise interaction is a cyclic convolution of \(\eta\)
spectra (in the continuum limit \(\delta Q\) =
R\(\oint\)\(\partial\)\_\(\eta\)(s\(^2\)) = 0), and conservation of the
integer charge Q\_wind is merely a corollary for in-band states (exact
conservation CV \(\approx\) 0 measured in the charged census
construction, Figure P6a). The apparent violation is the fold-back of
the doubling walk reaching Nyquist (37\% edge-power accumulation;
correlation -0.93, Figure P6b). On the finite register the doubling
ladder neutralizes in 4 steps: 1 \(\to\) 2 \(\to\) 4 \(\to\) 8 \(\to\)
0.

\begin{figure}
\centering
\pandocbounded{\includegraphics[keepaspectratio,alt={Figure P6}]{figs/fig_p6_cyclic_conservation_v1.png}}
\caption{Figure P6}
\end{figure}

\textbf{Figure P6}: Cyclic conservation and Nyquist fold-back.

\subsubsection{4.3 Readout rectification (readout-side
measurement)}\label{readout-rectification-readout-side-measurement}

Against all products of the sea-driven walk (the doubling-and-negation
ladder \(~2^b\cdot(\pm 1)\)), \textbf{only the observation clock J = 3
reads the entire charged content as \textbar q\textbar{} = 1}
(concentration 1.000 --- in the +1-seeded universe and in the +2-seeded
universe alike). J = 4 erases charge; J = 5, 6 scatter it into multiple
values (Figure P4). The algebraic basis is 2\^{}b mod 3 =
\(\pm\)\textasciitilde1. That the observation clock has denominator 3
rests on prior measurements {[}6{]}{[}7{]} --- the locking tongue of
denominator 3 dominates with a width ratio \textgreater{} 461. The value
of the elementary charge itself connects to the published measurement
{[}6{]} as the rational address sin\(^2\)(23\(\pi\)/124) =
\(\sqrt{4\pi\alpha}\) (agreement 16 parts in \(~10^{8}\)).

\begin{figure}
\centering
\pandocbounded{\includegraphics[keepaspectratio,alt={Figure P4}]{figs/fig_p4_rectification_v1.png}}
\caption{Figure P4}
\end{figure}

\textbf{Figure P4}: Readout rectification: only the denominator-3
observation clock reads all charged content as \textbar q\textbar{} = 1.

The bookkeeping also closes: in conserving constructions, \(\Delta Q_3\)
= -3\(\Delta W\) holds to a precision of 7\(\times\)\(~10^{-10}\) (Q3 =
net charge readable mod 3; W = the ledger of charge hidden in
mod-3-neutral composites). \textbf{Readable charge never disappears ---
its decrease equals, exactly, the transport into composites that the
denominator-3 clock reads as neutral} (Figure P5).

\begin{figure}
\centering
\pandocbounded{\includegraphics[keepaspectratio,alt={Figure P5}]{figs/fig_p5_ledger_v1.png}}
\caption{Figure P5}
\end{figure}

\textbf{Figure P5}: The bookkeeping identity \(\Delta Q_3\) =
-3\(\Delta W\).

\subsubsection{4.4 Hypothesis (Pillar 2): Q\_em = m/3 and
confinement}\label{hypothesis-pillar-2-q_em-m3-and-confinement}

We propose the charge rule that binds the above: \textbf{charge is Q\_em
= m/3, and the dominant observation clock (denominator 3) reads ``3 raw
windings = 1 elementary charge.''} Consequences:

\begin{itemize}
\tightlist
\item
  u-type m = +2 \(\to\) +2/3; d-type m = -1 \(\to\) -1/3; electron m =
  -3 \(\to\) -1; \(\nu\) m = 0 \(\to\) 0.
\item
  \textbf{Confinement = mod-3 readability}: the observer's space is the
  \(Z_3\) quotient of the \(\eta\) circle, and only species
  single-valued on it (m \(\equiv\) 0 mod 3) can be read out freely.
  Quarks (m \(\not\equiv\) 0) are unreadable alone --- \textbf{the 1/3
  charge and confinement are two faces of the same fact} and require no
  separate explanations.
\item
  Checks: proton uud = 2+2-1 = +3 \(\to\) Q = +1. Neutron udd = 2-1-1 =
  0 \(\to\) Q = 0. \(\pi\)\(^+\)(\(u\bar d\)) = 2+1 = +3 \(\to\) +1.
  \textbf{Integer charge for all hadrons follows automatically.}
\end{itemize}

\textbf{Confinement test (measured)}: quark-type (m = +2) + sea starts
at readable fraction f\_read = 0.0000 and becomes readable only to the
extent that the walk forms m \(\equiv\) 0 composites (hadron analogues)
(\(\to\) 0.2504 over 4000 collisions). Electron-type (m = +3) + sea
starts at f\_read = 1.0000 (free from the outset). \textbf{The isolated
quark-type system remains at f\_read = 0.0000 even after 4000
collisions} --- ``a single quark cannot be observed'' is realized as a
consequence of readout (Figure P9).

\begin{figure}
\centering
\pandocbounded{\includegraphics[keepaspectratio,alt={Figure P9}]{figs/fig_p9_confinement_v1.png}}
\caption{Figure P9}
\end{figure}

\textbf{Figure P9}: The confinement = mod-3 readability test.

Correspondence with the standard theory: this is the same structure as
the superselection by the SU(3) center \(Z_3\) --- ``only triality-0
states are free'' {[}10{]}{[}11{]}. But whereas in the standard theory
\(Z_3\) enters as an axiom on the dynamics side (the center of the gauge
group), in the present hypothesis mod 3 comes from the \textbf{readout
side}, as the measured dominant denominator of the observation clock ---
one of the novelties of this paper (§9).

\textbf{Two precisions}: (i) two logics in this section must be
distinguished --- ``the denominator-3 clock rectifies the charged walk
to \(\pm\)\textasciitilde1'' is a measurement; ``identifying the
physical charge as m/3'' is a hypothesis. To promote the identification
to a law requires a dynamical test of \textbf{response} to charge ---
whether the interaction strength \(\Gamma_{\rm int}\) of species with
different m scales as m/3 (or (m/3)\(^2\)) (§13). (ii) The confinement
of this paper is \textbf{readout-level confinement} (single-species
unreadability); the separation-energy growth, flux tubes, and asymptotic
freedom contained in standard confinement are not yet derived.
Precisely, it is a \textbf{readout precursor structure of confinement};
the test of whether the cost of separating two unreadable species grows
with distance (§13) is the promotion condition to dynamical confinement.

\subsection{\texorpdfstring{5. Pillar 3: the statistics rule --- the
\(Z_2\) clock
cover}{5. Pillar 3: the statistics rule --- the Z\_2 clock cover}}\label{pillar-3-the-statistics-rule-the-z_2-clock-cover}

\subsubsection{5.1 Statistics does not live in spatial rotation
(measured)}\label{statistics-does-not-live-in-spatial-rotation-measured}

The global rotation of the channel doublet is an exact symmetry of the
dynamics (all observables machine-zero flat in the control experiment;
the rotation invariance of Im(\(\bar b a\)) verified analytically), and
the spatial rotation charge is \(\ell\) = \(\pm\)\textasciitilde1 for
all observed modes (after calibration, deviations within
\(\pm\)\textasciitilde0.01 at N = 6, 8) --- neither half-integers nor
\(\ell\) = 2 appear among occupied modes. The two-valuedness of
statistics lives neither on the channel side nor on the spatial side.

\subsubsection{5.2 Statistics lives in the double cover of the clock
(measured)}\label{statistics-lives-in-the-double-cover-of-the-clock-measured}

Tracking the state phase of a species along the sequence of
observational recurrences (peaks of the autocorrelation of the
observable field s = Im(\(\bar b a\))), \textbf{the phase of the charged
species is exactly quantized to \{0, \(\pi\)\} at every recurrence
point} (\(\Phi\)/\(\pi\) = \(\pm\)\textasciitilde0.999 or +0.002;
3-digit precision; no continuous values) --- the state commutes between
the +1 and -1 sheets (Figure P8). The neutral non-projected bundle shows
a continuous phase; the presence or absence of quantization separates
the species.

\begin{figure}
\centering
\pandocbounded{\includegraphics[keepaspectratio,alt={Figure P8}]{figs/fig_p8_z2_quantization_v1.png}}
\caption{Figure P8}
\end{figure}

\textbf{Figure P8}: \(Z_2\) phase quantization: the charged species has
\(\Phi\) \(\in\) \{0, \(\pi\)\}; the neutral bundle is continuous.

The cross-table experiment (4 cells of \(\chi\) parity \(\times\)
winding \{1, 2\}) showed that \textbf{the covering degree depends only
on \(\chi\) parity, not on winding (charge)}: fermionic species (even
\(\chi\) frequency = odd harmonics) have covering degree 2 (\(Z_2\),
both sheets); bosonic species have degree 1. Furthermore, a pure m = 0
fermionic-band species (winding-shift construction; \(\eta\) occupancy
(0, 1.000)) also showed covering degree 2 with Qz2 = 1.00 ---
\textbf{the row of the neutral fermion (neutrino) is established}
(Figure P10).

\begin{figure}
\centering
\pandocbounded{\includegraphics[keepaspectratio,alt={Figure P10}]{figs/fig_p10_nu_spinstat_v1.png}}
\caption{Figure P10}
\end{figure}

\textbf{Figure P10}: Establishment of the \(\nu\) row and the
spin-statistics correspondence.

\subsubsection{5.3 Hypothesis (Pillar 3) and prior
contrast}\label{hypothesis-pillar-3-and-prior-contrast}

We propose \textbf{Fermi/Bose distinction = clock covering degree} (1 =
bosonic / 2 = fermionic). This is the counterpart of the structure by
which Finkelstein--Rubinstein {[}12{]} derived soliton spin-statistics
from the double cover of configuration space --- but here \textbf{the
double cover lives not in space (rotation/exchange) but in the clock
(state period = 2 \(\times\) observation period)}. The same structure
connects to the previously measured fidelity of order 248 / observation
124 {[}6{]}{[}7{]} (recurrence invisible at \(F_{124}\); \(F_{248}\) =
1).

\subsection{6. Pillar 4: the mass rule --- mass is a relation to the
sea}\label{pillar-4-the-mass-rule-mass-is-a-relation-to-the-sea}

\subsubsection{6.1 Equivalence theorems (three,
measured)}\label{equivalence-theorems-three-measured}

\begin{itemize}
\tightlist
\item
  \textbf{Winding-shift equivalence}: e\^{}\{im\(\eta\)\} is an exact
  symmetry of the dynamics (\(\bar b a\) invariant); isolated pure
  address species are all unitarily equivalent --- measured: mass\(^2\),
  polarization, and s\_z agree to machine precision across all
  addresses. \textbf{Mass and spin are not attributes of the isolated
  address.}
\item
  \textbf{Divisor-class theorem}: in the sea, properties differentiate
  --- but \textbf{only down to the divisor class gcd(m, 16)} (mass\(^2\)
  of the odd addresses \{1,3,5,7\} agrees to 5 digits; \{2,6\} agree;
  \{4\} is distinct; Figure P7). The reason is exact: the unit
  automorphisms m \(\to\) um of \(Z_{16}\) fix the sea (m = 0), so sea
  coupling cannot distinguish unit addresses in principle. Consequence:
  \textbf{the distinction between \(\pm\)\textasciitilde1 and
  \(\pm\)\textasciitilde5, \(\pm\)\textasciitilde7 does not exist in the
  dynamics; only the mod-3 readout breaks the unit class} --- an
  independent corroboration of Pillar 2 (readout rectification).
\item
  \textbf{\(\chi\)-band translation symmetry}: placing the same-winding
  species in \(\chi\) bands (10,12,14)/(30,32,34)/(50,52,54) leaves
  mass\(^2\) identical to 5 digits --- band position does not
  differentiate species either (§10, weakness 1).
\end{itemize}

\begin{figure}
\centering
\pandocbounded{\includegraphics[keepaspectratio,alt={Figure P7}]{figs/fig_p7_divisor_class_v1.png}}
\caption{Figure P7}
\end{figure}

\textbf{Figure P7}: The divisor-class theorem: in-sea properties depend
only on gcd(m, 16).

\subsubsection{6.2 Hypothesis (Pillar 4)}\label{hypothesis-pillar-4}

\textbf{Mass is a relational quantity with the sea (vacuum).} Only the
species locked to the clock (\(\rho\) = 1/1) is massless and stable (the
photon-like ground species; mass degree \(~10^{-11}\) measured); content
away from the clock carries mass as resonances, and the sea
differentiates the individuality of species (lifetime, polarization,
mass) at the level of divisor classes. The view of the vacuum-as-sea as
co-determinant of particle properties resonates with the lineage of
emergent gauge fields and emergent fermions {[}13{]} (string-net
condensation; superfluid \(^3\)He), but this paper differs in having
measured the sea dependence of properties in the form of equivalence
theorems.

\subsection{7. Pillar 5: the periodic table of waves --- double
signboard and the 62-species
assignment}\label{pillar-5-the-periodic-table-of-waves-double-signboard-and-the-62-species-assignment}

\subsubsection{7.1 The double signboard}\label{the-double-signboard}

\begin{itemize}
\tightlist
\item
  \textbf{Table A (isolated stable-element table)}: pure address species
  on the U\^{}144 clock lattice, all stable by self-replication
  (measured). Algebraic classifiers: readable charge (mod 3), doubling
  orbit, neutralization steps, \(\eta\) parity. A second parity
  structure: odd addresses resist doubling-neutralization maximally (4
  steps), even ones fall fast (\(\pm\)\textasciitilde4 in 2 steps, +8 in
  1).
\item
  \textbf{Table B (in-sea lifetime table)}: sea coupling selects species
  --- charged long-lived (\(\tau\) \textasciitilde{} \(~10^{4}\)),
  neutral transients short-lived (\(\tau\) \textasciitilde{}
  \(~10^{2\text{--}3}\)), composites neutralized. \textbf{Observed
  ``particles'' are this quasi-stable hierarchy.}
\end{itemize}

\textbf{The native periodic table (model-specific, measurement-based
master table)}: we first write the table using only this paper's
classifiers, without Standard-Model particle names. The Standard-Model
assignment (§7.2) is a hypothesis on top of this table.

{\def\LTcaptype{none} % do not increment counter
\begin{longtable}[]{@{}
  >{\raggedright\arraybackslash}p{(\linewidth - 12\tabcolsep) * \real{0.1429}}
  >{\raggedright\arraybackslash}p{(\linewidth - 12\tabcolsep) * \real{0.1429}}
  >{\raggedright\arraybackslash}p{(\linewidth - 12\tabcolsep) * \real{0.1429}}
  >{\raggedright\arraybackslash}p{(\linewidth - 12\tabcolsep) * \real{0.1429}}
  >{\raggedright\arraybackslash}p{(\linewidth - 12\tabcolsep) * \real{0.1429}}
  >{\raggedright\arraybackslash}p{(\linewidth - 12\tabcolsep) * \real{0.1429}}
  >{\raggedright\arraybackslash}p{(\linewidth - 12\tabcolsep) * \real{0.1429}}@{}}
\toprule\noalign{}
\begin{minipage}[b]{\linewidth}\raggedright
Raw winding class
\end{minipage} & \begin{minipage}[b]{\linewidth}\raggedright
mod-3 readable
\end{minipage} & \begin{minipage}[b]{\linewidth}\raggedright
Neutralization steps
\end{minipage} & \begin{minipage}[b]{\linewidth}\raggedright
Clock cover (F-band/B-band)
\end{minipage} & \begin{minipage}[b]{\linewidth}\raggedright
In-sea mass\(^2\) (divisor class)
\end{minipage} & \begin{minipage}[b]{\linewidth}\raggedright
Isolated stability
\end{minipage} & \begin{minipage}[b]{\linewidth}\raggedright
In-sea lifetime
\end{minipage} \\
\midrule\noalign{}
\endhead
\bottomrule\noalign{}
\endlastfoot
m = 0 & 0 (free, neutral) & 0 & 2 / 1 & --- (sea, ground species) &
stable & ground species stable \\
odd (units)
\{\(\pm\)\textsubscript{1,\(\pm\)}3,\(\pm\)\textsubscript{5,\(\pm\)}7\}
& \(\pm\)\textasciitilde1 or 0 & 4 (max) & 2 / 1 & 0.787 (5-digit
agreement) & exactly stable & \(\tau\) \textasciitilde{} \(~10^{4}\)
(charged) \\
\{\(\pm\)\textsubscript{2,\(\pm\)}6\} & \(\mp\)\textasciitilde1 or 0 & 3
& 2 / 1 & 0.764 & exactly stable & retention 0.226 \\
\{\(\pm\)\textasciitilde4\} & \(\pm\)\textasciitilde1 & 2 & 2 / 1 &
0.722 & exactly stable & retention 0.354 (longest) \\
\{+8\} & -1 & 1 (shortest) & --- & unmeasured & --- & unmeasured \\
\end{longtable}
}

The covering degree is determined by \(\chi\) parity, not by the raw
winding class (cross-table measurement), hence an independent column.

\subsubsection{7.2 The 62-species Standard-Model
assignment}\label{the-62-species-standard-model-assignment}

Confidence tags: \textbf{E} = measurement anchor (directly supported by
measurements in this or published papers) / \textbf{S} = structural
correspondence (mechanism matches; quantitatively unverified) /
\textbf{H} = hypothesis (future verification target).

{\def\LTcaptype{none} % do not increment counter
\begin{longtable}[]{@{}
  >{\raggedright\arraybackslash}p{(\linewidth - 10\tabcolsep) * \real{0.1667}}
  >{\raggedright\arraybackslash}p{(\linewidth - 10\tabcolsep) * \real{0.1667}}
  >{\raggedright\arraybackslash}p{(\linewidth - 10\tabcolsep) * \real{0.1667}}
  >{\raggedright\arraybackslash}p{(\linewidth - 10\tabcolsep) * \real{0.1667}}
  >{\raggedright\arraybackslash}p{(\linewidth - 10\tabcolsep) * \real{0.1667}}
  >{\raggedright\arraybackslash}p{(\linewidth - 10\tabcolsep) * \real{0.1667}}@{}}
\toprule\noalign{}
\begin{minipage}[b]{\linewidth}\raggedright
Particle
\end{minipage} & \begin{minipage}[b]{\linewidth}\raggedright
States
\end{minipage} & \begin{minipage}[b]{\linewidth}\raggedright
Winding m (Q\_em = m/3)
\end{minipage} & \begin{minipage}[b]{\linewidth}\raggedright
Cover
\end{minipage} & \begin{minipage}[b]{\linewidth}\raggedright
Conf.
\end{minipage} & \begin{minipage}[b]{\linewidth}\raggedright
Remarks
\end{minipage} \\
\midrule\noalign{}
\endhead
\bottomrule\noalign{}
\endlastfoot
u, c, t & 18 & +2 (+2/3) & 2 & \textbf{E} & m \(\not\equiv\) 0 \(\to\)
confinement (Fig. P9 measured). Color = mod-3-unreadable raw-winding
residue (ledger W) \\
d, s, b & 18 & -1 (-1/3) & 2 & \textbf{E} & same \\
e, \(\mu\), \(\tau\) & 6 & -3 (-1) & 2 & \textbf{E} & free (m \(\equiv\)
0). Elementary-charge address sin\(^2\)(23\(\pi\)/124) {[}6{]} \\
\(\nu\) \(\times\)\textasciitilde3 & 6 & 0 (0) & 2 & \textbf{E} &
neutral fermion (established by measurement, §5.2) \\
\(\gamma\) & 1 & 0 & 1 & \textbf{E} & ground species \(\rho\) = 1/1
(massless, stable, measured) \\
g & 8 & color pairs (mod-3 neutral, raw winding nonzero) & 1 & S & color
non-singlet \(\to\) unreadable alone = confinement. The octet's
generation not yet derived \\
W\(\pm\), Z & 3 & \(\pm\)\textasciitilde3, 0 & 1 & S & t-winding
(own-clock detuning) \(\to\) massive (consistent with entrainment
measurements; dynamics unverified) \\
H & 1 & 0 & 1 & S & collective mode of the sea (condensate). \(\ell\) =
0 \\
G & 1 & 0 & 1 & H & only as a composite of two \(\ell\) = 1 quanta (§8.4
prediction). \(\ell\) = \(\pm\)\textasciitilde2 \\
(generation axis) & --- & --- & --- & H & the axis differentiating 3
generations is unidentified (§10) \\
\end{longtable}
}

Total \textbf{62 species} (36 quarks + 12 leptons + 8 gluons +
\(\gamma\) + 2 W + Z + H + G). Antiparticles are m \(\to\) -m (charge
conjugation = \(\eta\) reversal; pair production measured {[}3{]}). The
generation axis g = 1, 2, 3 is unresolved (§10). Per-row confidence is
spelled out in the attached periodic table v0.4.

\subsubsection{7.3 The relation between the two finite structures (open
structural
problem)}\label{the-relation-between-the-two-finite-structures-open-structural-problem}

Two finite structures are measured in this paper: the universal clock
U\^{}144 (§3) and the winding register \(Z_{16}\) (§4.2). The factor
structure is 144 = \(~2^4\cdot 3^2\), 16 = \(~2^4\), 144/16 = 9 =
\(~3^2\), and this paper's classifiers correspond to the factors ---
\textbf{2 = statistics cover (§5), 3 = charge readability (§4), 16 = raw
winding cycle (§4.2)}. Whether this coincidence is accidental, or a
resonance structure in which both follow from the update rule (a
resonance classification of clock order \(\times\) internal winding
order), is the largest open structural problem of this paper, registered
in §13 as a task to be tested \textbf{as an operator recurrence
structure}, not as numerology. Formalization: bring every readout
operation into the same operator representation and measure the order
ord(U\textbar\_\{H\_i\}) per subspace --- the true periodic law may be
the family of ``orders that a single operator has on different readout
subspaces.'' The true periodic law of the ``periodic table'' probably
lives here.

\subsection{8. Explanatory power --- connections to existing
problems}\label{explanatory-power-connections-to-existing-problems}

\subsubsection{\texorpdfstring{8.1 Resolution of the previous paper's
open problem, ``charge does not stabilize at
\(\pm\)\textasciitilde1''}{8.1 Resolution of the previous paper's open problem, ``charge does not stabilize at \textbackslash pm\textasciitilde1''}}\label{resolution-of-the-previous-papers-open-problem-charge-does-not-stabilize-at-pm1}

The previous paper {[}5{]} §9.5 recorded as an open problem that
``charge (winding) spreads over multiple quantized values and does not
stabilize at \(\pm\)\textasciitilde1.'' Under the present hypothesis
this is no contradiction: on the \textbf{state side}, the sea-driven
walk (measured, §4.1) keeps dispersing windings. On the \textbf{readout
side}, the denominator-3 observation clock folds all its products back
to \textbar q\textbar{} = 1 (measured, §4.3). ``Multivaluedness of
charge'' and ``universality of the elementary charge'' coexist without
contradiction --- an open problem turns into a prediction of the
hypothesis.

\subsubsection{8.2 The 1/3-charge problem of
quarks}\label{the-13-charge-problem-of-quarks}

Under Q\_em = m/3, fractional charge (why 1/3?) and confinement (why
never alone?) are a single mechanism --- mod-3 readability (§4.4;
measured, Fig. P9).

\subsubsection{8.3 Antimatter}\label{antimatter}

Charge conjugation m \(\to\) -m is a symmetry of the spectrum, and pair
production is a necessity of the walk (published measurement {[}3{]}).

\subsubsection{8.4 The weakness of gravity (a bold prediction ---
classified as a prediction, not a
pillar)}\label{the-weakness-of-gravity-a-bold-prediction-classified-as-a-prediction-not-a-pillar}

The measurements are two: the rotation-generator spectrum of the
condensate has no linear \(\ell\) = 2 slot (calibrated ratio at most
1.000, exact); and the spectrum of the squared readout carries a
coherent 2\(\omega\) quadrupole line at all N (ratio 2 \(\pm\) 6\%). The
strong hypothesis that follows: \textbf{if a gravitational spin-2 mode
exists, it appears not as a first-order degree of freedom but only as a
second-order composite of two \(\ell\) = 1 quanta.} Beyond that --- the
quadratic suppression of the coupling (consistent with the prior
derivation \(\alpha_G\) = (m/M)\(^2\) {[}15{]}), universal attraction,
and the isomorphism with the double copy {[}16{]} --- are registered as
\textbf{predictions} for future dynamical verification. If the
prediction is right, the reason gravity is weak and the reason it is
hard to verify are one and the same structure.

\subsection{9. Relation to prior work, and the delimitation of
novelty}\label{relation-to-prior-work-and-the-delimitation-of-novelty}

Each component of the hypothesis has a strong prior lineage: winding =
charge (Skyrme {[}17{]}; Kaluza--Klein {[}18{]}); mod-3 confinement
(triality {[}10{]}{[}11{]}); the double cover of statistics
(Finkelstein--Rubinstein {[}12{]}); vacuum = sea (Wen; Volovik
{[}13{]}); gravity = square (KLT/BCJ {[}16{]}); full-assignment attempts
(the Eightfold Way {[}9{]}; preons {[}14{]}). \textbf{This paper claims
no novelty for the components.} The novelty is limited to three points:

\begin{enumerate}
\def\labelenumi{\arabic{enumi}.}
\tightlist
\item
  That these components come out \textbf{simultaneously from
  measurements of a single anonymous dynamics} that assumes no particle
  concept (confluence).
\item
  The presentation of a counterpart in which charge unit, confinement,
  and elementary-charge uniqueness come not from a dynamics-side axiom
  (the center of a gauge group) but from the \textbf{divisor structure
  of the observation clock (the readout side)}.
\item
  The measurement that the double cover of statistics lives not in space
  (rotation of configuration space) but in the \textbf{clock} (state
  period = 2 \(\times\) observation period).
\end{enumerate}

\subsection{10. Weaknesses and unverified parts (honest
registration)}\label{weaknesses-and-unverified-parts-honest-registration}

\begin{enumerate}
\def\labelenumi{\arabic{enumi}.}
\tightlist
\item
  \textbf{The origin of generations is unresolved}: the naive hypothesis
  generation = \(\chi\)-band position was refuted (\(\chi\)-band
  translation symmetry, §6.1). Band ratios, crystallographic orders
  \{3,4,6\}, and composite hierarchy remain candidates.
\item
  \textbf{The W/Z rows}: direct production is impossible in the
  low-energy laboratory (N \(\le\) 16, small sea) --- O(1) detuning
  decays at the register edge (structural). The \(~1/\delta^2\) scaling
  test of the effective contact vertex (the Fermi-theory road) is
  specified but not yet executed.
\item
  \textbf{The graviton row}: the quadrupole 2\(\omega\) line exists at
  all N, but the pre-registered criterion (\textbar ratio-2\textbar{}
  \textless{} 0.02) passed only at 1 of 3 --- frame calibration is
  incomplete.
\item
  \textbf{Dependence on the register order ne = 16}: the details of
  Table A (neutralization steps etc.) are functions of ne.
  Re-examination with variable ne is needed (also a testable prediction:
  ``the period of the periodic table is a function of the register
  order'').
\item
  \textbf{The unification of Q\_em = m/3 with the transmission charge
  \(\sqrt{4\pi\alpha}\) {[}6{]}{[}19{]} is incomplete} --- the relation
  between the ``unit'' (3 windings) and the ``value'' (0.302822) of the
  elementary charge is a task for the next paper.
\item
  Part of the sea constructions (artificial projected seas) was found to
  lie outside the conservation class (Nyquist fold-back, §4.2). The
  corresponding quantitative results are treated with reservation;
  qualitative conclusions were cross-checked in conserving-class series.
\end{enumerate}

\subsection{11. Falsification
conditions}\label{falsification-conditions}

\begin{enumerate}
\def\labelenumi{\arabic{enumi}.}
\tightlist
\item
  If an observation clock with a denominator other than 3 is constructed
  that rectifies the entire set of charged-walk products to a unique
  elementary-charge unit, Pillar 2 is rejected.
\item
  If, in the same condensed phase, with the same update rule and the
  same settling criterion, a series is obtained whose \(\omega\)
  deviates systematically from \(\pi\)/72 in the long-window limit,
  Pillar 1 is rejected.
\item
  If a species is constructed whose covering degree varies independently
  of \(\chi\) parity, Pillar 3 is rejected.
\item
  If a measurement shows species properties (mass, lifetime,
  polarization) differentiating without a sea, Pillar 4 is rejected.
\item
  If a construction is found in which an m \(\not\equiv\) 0 (mod 3)
  species becomes readable alone, the confinement rule is rejected.
\end{enumerate}

\subsection{12. Record of refutations and corrections (13 items,
summarized)}\label{record-of-refutations-and-corrections-13-items-summarized}

All hypotheses and instruments refuted in the course of this work are
recorded: (1) short-window sideband addresses are transients, not stable
species (vanish in long windows). (2) The mass = detuning\(^2\) law was
an apparent relation on transients. (3) The naive form of
\(\pm\)\textasciitilde1 fixed-point uniqueness (every pure m is a fixed
point; uniqueness does not follow). (4) No mass-width correlation (r =
-0.07). (5) The sea-construction convention (\(\chi\)-analytic
projection) hypothesis failed --- the true cause of conservation
violation is Nyquist fold-back (§4.2). (6) \(\pm\)\textasciitilde1
uniqueness by escape rates is undecidable in principle, by the
divisor-class theorem. (7) The kinematic version of SU(2) recurrence
discrimination is spinorial for all states (trivial). (8) Single-point
covering judgments pick up accidental proximity to \(\pi\) (calibrated
to recurrence-sequence judgment). (9) No linear \(\ell\) = 2 slot for
the graviton (reinterpreted as composite). (10) Generation =
\(\chi\)-band position refuted. (11) The band-shift method of W/Z
detuning injection failed (effectively linear dispersion). (12) The
projection method of \(\nu\) construction failed (bundle hair m = +1 was
the cause; solved by the shift method). (13) The earlier projected seas
were amplified numerical residue (legitimate as m = 0 fields;
qualitative results unaffected; recorded). Details are in the attached
analysis notes.

\subsection{13. Requirements handed to the unified dynamics
(specification for the next
paper)}\label{requirements-handed-to-the-unified-dynamics-specification-for-the-next-paper}

To verify this classification hypothesis at the dynamics level, the two
currently separate dynamics must be unified into a single
\textbf{universal interaction function}, and the 62 species must be
shown to be generated, to decay, and to interact as its spectrum. The
measurements of this paper fix the requirements that the function must
satisfy:

\begin{enumerate}
\def\labelenumi{\arabic{enumi}.}
\tightlist
\item
  \textbf{Anonymity}: wave forms are the only input; no special-casing
  by species (no IF statements).
\item
  \textbf{All channels simultaneously}: electromagnetic (winding
  exchange), weak (clock detuning), strong (mod-3-unreadable raw-winding
  residue), and gravitational (second-moment coupling) must come out not
  as separate terms but as different readouts of the same operation,
  driving the three planes (xy = spin / zR = mass / tQ = charge)
  {[}20{]} simultaneously.
\item
  \textbf{Automatic conservation}: cyclic winding conservation (§4.2),
  closure conservation, and norm conservation must hold from the
  structure of the operation, not by design.
\item
  \textbf{Reproduction of the universal clock}: the U\^{}144 clock
  lattice with \(\omega\) = \(\pi\)/72 (§3) must appear as spectrum.
\item
  \textbf{Emergent selection rules}: the sum rule m* = 2m\_B - m\_s, the
  doubling-and-negation ladder, \(Z_2\) phase quantization, and the
  divisor-class structure must come out without hand placement.
\item
  \textbf{Exact computability}: preserve anonymity \(\Rightarrow\)
  computation linear in state size {[}4{]} (closed form or
  integer-exact).
\item
  \textbf{Specialization to the two systems}: the two-body closed-form
  solution {[}3{]} and the N-body relational-wave dynamics {[}2{]} must
  be re-derived as restrictions/limits of the function.
\item
  \textbf{Falsifiable form}: if any of the 62 assigned species (§7.2)
  fails to appear in the spectrum, the periodic-table hypothesis is
  rejected row by row.
\end{enumerate}

\subsubsection{13.2 Verification program (top
four)}\label{verification-program-top-four}

In parallel with constructing the universal function, we register, in
order of priority, experiments that can directly strengthen or refute
this paper's hypothesis:

\begin{enumerate}
\def\labelenumi{\arabic{enumi}.}
\tightlist
\item
  \textbf{Dynamical derivation of denominator 3, and charge response}:
  the dominance of denominator 3 currently rests on locking measurements
  {[}20{]}. If an analytic derivation is obtained of why the
  denominator-3 Arnold tongue is maximal in this dynamics, the ``3'' of
  Q\_em = m/3 advances from an empirical selection to a dynamical
  necessity. In addition, measure whether the exchange quantities and
  phase shifts of two species (m\(_1\), m\(_2\)) scale as the product
  m\(_1\)m\(_2\)/9 (a Coulomb-type charge-product test) --- if it
  passes, m/3 is promoted from bookkeeping to the physical charge
  appearing in interactions.
\item
  \textbf{Clock covering degree and exchange sign in a single experiment
  (the decisive experiment to be performed before constructing the
  universal function)}: connect covering degree 2
  \(\Longleftrightarrow\) exchange sign -1 of two identical species in
  one experiment. If the exchange operation E and the
  one-clock-revolution operation T are shown to belong to the same
  nontrivial class of the same double cover (E \(\simeq\) T, E\(^2\) =
  T\(^2\) = I), then half-integer recurrence, exchange sign, exclusion,
  and clock covering degree contract to a single \(Z_2\), and Pillar 3
  advances from ``resembling statistics'' to ``a reconstruction of the
  structure of statistics.''
\item
  \textbf{Deformation law of the table under variable register order}:
  with ne = 12, 18, 24, 32, \ldots, measure whether divisor classes,
  neutralization steps, readability, and covering structure follow
  gcd(m, ne)/Div(ne). If they follow, the divisor-class theorem
  generalizes; if not, 16 hides a structure of its own.
\item
  \textbf{Blind predictions}: without using Standard-Model particle
  names, predict from the model alone the unoccupied classification
  cells, lifetime rankings, charges, covering degrees, and decay
  destinations, and only afterwards compare with the known particle
  table --- a test format with the \(\Omega^-\)-type decisive power of
  the Eightfold Way.
\end{enumerate}

\subsection{14. Reproducibility}\label{reproducibility}

All experiments are published as deterministic scripts (28) in the same
folder as this paper. The correspondence table (the experiment inventory
file (jikken\_ichiran\_v1.md)) gives the complete mapping of scripts
\(\leftrightarrow\) result JSONs (25) \(\leftrightarrow\) figures (10)
\(\leftrightarrow\) claims. The three items originally executed
interactively were formalized as scripts, and reproduction agreement
with the saved results was machine-verified. Judgment criteria were
fixed and recorded in each script header before execution.

\subsection{References}\label{references}

\textbf{Self-citations (sources of dynamics and measurements)} {[}1{]}
N. Kihara, Thought Experiment ``On the R Axis,'' Zenodo (2026).
doi:10.5281/zenodo.19902677 {[}2{]} N. Kihara, Causal Separation of the
Metastable Phase by Two-Stage Seed Removal (Paper 8), Zenodo (2026).
doi:10.5281/zenodo.21614402 {[}3{]} N. Kihara, Generation of Fermionic
Structure --- the Universal Inelastic Map, Zenodo (2026).
doi:10.5281/zenodo.21808091 {[}4{]} N. Kihara, How Inflation Ends Is How
Matter Generation Begins (N-Body Connection), Zenodo (2026).
doi:10.5281/zenodo.21809814 {[}5{]} N. Kihara, Genesis of the Three
Spatial Axes and Proper Time, Zenodo (2026). doi:10.5281/zenodo.21816651
{[}6{]} N. Kihara, Counting Readout in an Anonymous Two-Channel Closed
Wave System (Trilogy B), Zenodo (2026). doi:10.5281/zenodo.21763997
{[}7{]} N. Kihara, The Generation Structure of Fermions (Paper 9),
Zenodo (2026). doi:10.5281/zenodo.21766706 {[}15{]} N. Kihara,
Re-Reading Centered on Reality --- the Derivation-Replacement Edition,
Zenodo (2026). doi:10.5281/zenodo.21765367 {[}19{]} N. Kihara,
Two-Grammar Decomposition in an Anonymous Two-Channel Closed Wave System
(Trilogy A), Zenodo (2026). doi:10.5281/zenodo.21763995 {[}20{]} N.
Kihara, Locking Dynamics in an Anonymous Two-Channel Closed Wave System
(Trilogy C), Zenodo (2026). doi:10.5281/zenodo.21763999

\textbf{External references} {[}8{]} W. Thomson (Lord Kelvin), ``On
Vortex Atoms,'' Proc. Roy. Soc. Edinburgh 6, 94 (1867). {[}9{]} M.
Gell-Mann, ``The Eightfold Way,'' Caltech Report CTSL-20 (1961); M.
Gell-Mann and Y. Ne'eman, The Eightfold Way (Benjamin, 1964). {[}10{]}
K. G. Wilson, ``Confinement of quarks,'' Phys. Rev.~D 10, 2445 (1974).
{[}11{]} G. 't Hooft, ``On the phase transition towards permanent quark
confinement,'' Nucl. Phys. B 138, 1 (1978). {[}12{]} D. Finkelstein and
J. Rubinstein, ``Connection between spin, statistics, and kinks,'' J.
Math. Phys. 9, 1762 (1968). {[}13{]} M. A. Levin and X.-G. Wen,
``String-net condensation,'' Phys. Rev.~B 71, 045110 (2005); G. E.
Volovik, The Universe in a Helium Droplet (Oxford, 2003). {[}14{]} H.
Harari, ``A schematic model of quarks and leptons,'' Phys. Lett. B 86,
83 (1979); M. A. Shupe, Phys. Lett. B 86, 87 (1979). {[}16{]} Z. Bern,
J. J. M. Carrasco, and H. Johansson, ``Perturbative quantum gravity as a
double copy of gauge theory,'' Phys. Rev.~Lett. 105, 061602 (2010); H.
Kawai, D. C. Lewellen, and S. H. H. Tye, Nucl. Phys. B 269, 1 (1986).
{[}17{]} T. H. R. Skyrme, ``A unified field theory of mesons and
baryons,'' Nucl. Phys. 31, 556 (1962). {[}18{]} T. Kaluza, Sitzungsber.
Preuss. Akad. Wiss. (1921) 966; O. Klein, Z. Phys. 37, 895 (1926). Note:
S. Coleman, ``Quantum sine-Gordon equation as the massive Thirring
model,'' Phys. Rev.~D 11, 2088 (1975), representative of the
bosonization lineage, is referenced via {[}7{]}.

\end{document}
